\section{A quantitative height for nonreal heat zeros}
\label{sec:height}
The following bound makes the existence of a nonreal zero at each small
negative time quantitative. Its analytic input is Dobner's uniform
approximation to the deformed Dirichlet series; the recurrence estimate
needed to place a zero at a finite height is proved here.

\begin{theorem}\label{thm:nonreal-height}
There are absolute constants $C,a_0>0$ such that, for $0<a<a_0$,
$H_{-a}$ has a simple zero $z_a$ satisfying
\[
 0<\Re z_a\le \exp\!\bigl(\exp(C/\sqrt a)\bigr),
 \qquad \Im z_a\le-\tfrac12.
\]
The complex conjugate of $z_a$ is another simple zero.
\end{theorem}

\begin{proof}
We use the normalization
\[
 \Xi_{-a}(s)=8H_{-a}(-i(2s-1)),\quad
 \gamma(s)=\tfrac12s(s-1)\pi^{-s/2}\Gamma(s/2),
 \quad \Xi_0(s)=\gamma(s)\zeta(s).
\]
Put
\[
 F_a(s)=\sum_{n\ge1}e^{-a(\log n)^2/4}n^{-s},\qquad
 J_a(s)=s+\frac a4\operatorname{Log}\frac{s}{2\pi},
 \quad
 \gamma_a(s)=\gamma(s)\exp\frac{(s-J_a(s))^2}{a}.
\]
The logarithm is principal in the upper half-plane. Dobner
\cite[Theorem~4, equations~(10)--(13)]{dobner} proves that absolute
constants $C_D,Y_D$ give
\begin{equation}\label{eq:dobner-height}
 \frac{\Xi_{-a}(J_a(s))}{\gamma_a(s)}
 =F_a(s)+O\bigl(C_D(\Im s)^{-1/5}e^{40/a}\bigr)
\end{equation}
uniformly for $0<a\le1$, $1/2\le\Re s\le2$ and $\Im s\ge Y_D$.
In particular, the threshold and error constant are independent of $a$.

The pole of the original zeta function produces the scaling function
\[
 E(w)=\int_0^\infty e^{-v^2/4-wv}\,dv.
\]
This function is entire and satisfies $|E(w)|\le2\sqrt\pi e^{|w|^2}$.
It has a zero: otherwise its entire logarithm $g$ has real part bounded
above by $O(1+|w|^2)$; Borel--Carath\'eodory and Cauchy's estimate imply
that $g$ is a polynomial of degree at most two. But
\[
 xE(x)=\int_0^\infty e^{-u-u^2/(4x^2)}\,du\longrightarrow1
 \quad(x\longrightarrow+\infty),
\]
so $\Re g(x)=-\log x+o(1)$, a contradiction. The integral is positive
on the real axis. Integration by parts gives $E'=2wE-2$, and thus every
zero of $E$ is nonreal and simple.

We shall also need an error estimate in this scaling. For
$s=1+\sqrt a\,w$, comparison of a sum with its integral yields
\begin{equation}\label{eq:heat-pole-profile}
 F_a(1+\sqrt a\,w)=a^{-1/2}E(w)+\mathcal E_a(w),
 \qquad |\mathcal E_a(w)|\le7,
\end{equation}
provided $\Re s\ge1/2$, $|s|\le2$ and $0<a\le1$. Indeed, for
$f_s(u)=e^{-a(\log u)^2/4}u^{-s}$ the error is at most
\[
 |f_s(1)|+\int_1^\infty|f_s'(u)|\,du
 \le 1+2\int_1^\infty u^{-3/2}du
       +\tfrac a2\int_1^\infty(\log u)u^{-3/2}du\le7.
\]
The substitution $v=\sqrt a\log u$ evaluates the integral exactly.

Fix a closed disk $D$ containing just one zero of $E$, and no real
point. Let $\mu=\min_{\partial D}|E|>0$, and choose $K\ge1$ with
$|w|\le K$ on $D$. Define
\[
 W_K=1+\int_0^\infty e^{-v^2/4+Kv}\,dv.
\]
Choose fixed $R>0$ and an integer $Q\ge2$ so that
\[
 \int_R^\infty e^{-v^2/4+Kv}\,dv\le\mu/64,
 \qquad 2\pi/Q\le\mu/(32W_K).
\]
For small $a$, put $N=\lceil e^{R/\sqrt a}\rceil$. The positive
majorant $e^{-a(\log u)^2/4}u^{-1+K\sqrt a}$ decreases for $u\ge1$;
therefore
\begin{align*}
 \sum_{n\ge1}e^{-a(\log n)^2/4}n^{-1+K\sqrt a}&\le W_K/\sqrt a,\\
 \sum_{n>N}e^{-a(\log n)^2/4}n^{-1+K\sqrt a}&\le\mu/(64\sqrt a).
\end{align*}
For any $Y>0$, pigeonhole the $Q^{N-1}+1$ points
\[
 \left(\left\{\frac{jY\log n}{2\pi}\right\}\right)_{2\le n\le N},
 \qquad 0\le j\le Q^{N-1},
\]
into $Q^{N-1}$ cubes of side $1/Q$. Their difference supplies
$\tau=kY$, $1\le k\le Q^{N-1}$, with
$|n^{-i\tau}-1|\le2\pi/Q$ for every $n\le N$. The last two estimates
then give
\begin{equation}\label{eq:heat-finite-recurrence}
 \sup_{w\in D}|F_a(1+\sqrt a\,w+i\tau)-F_a(1+\sqrt a\,w)|
 \le\frac{\mu}{16\sqrt a},\qquad
 Y\le\tau\le YQ^{N-1}.
\end{equation}

Take $Y=e^{250/a}$ and this value of $\tau$, and set
$s(w)=1+\sqrt a\,w+i\tau$. For small enough $a$, throughout $D$,
$3/4\le\Re s(w)\le5/4$ and $\Im s(w)\ge Y/2\ge Y_D$.
The function
\[
 B_a(w)=\Xi_{-a}(J_a(s(w)))/\gamma_a(s(w))
\]
is holomorphic near $D$. Equation~\eqref{eq:dobner-height} has error
at most $2^{1/5}C_De^{-10/a}\le\mu/(16\sqrt a)$. Combined with
\eqref{eq:heat-pole-profile} and \eqref{eq:heat-finite-recurrence}, this
shows on $\partial D$ that
\[
 |B_a-a^{-1/2}E|\le\frac{\mu}{8\sqrt a}+7
 <\frac{\mu}{\sqrt a}=|a^{-1/2}E|.
\]
Rouch\'e's theorem gives exactly one zero $w_a$ in $D$, with
multiplicity one. Write $s_a=s(w_a)$ and
$z_a=-i(2J_a(s_a)-1)$. Since $J_a'(s_a)=1+a/(4s_a)\ne0$,
the zero of $H_{-a}$ at $z_a$ is simple.

For small $a$, $|s_a|>2\pi$; hence
$\Re J_a(s_a)\ge3/4$ and $\Im z_a\le-1/2$.
Also $0<\Re z_a=2\Im J_a(s_a)\le3\tau$. Consequently
\[
 \log\Re z_a\le\log3+250/a+2(\log Q)e^{R/\sqrt a}
 \le e^{C/\sqrt a}
\]
for a fixed $C>R$ and all small enough $a$. Reality of $H_{-a}$ on
the real axis gives the conjugate zero.
\end{proof}
